% Extracted section; already included in the complete working note.
\section{The Lie-algebra interpretation: squeezing and imaginary time}
\label{sec:Lie-algebra}
The supplied preliminary notes \cite{Notes} provide a useful algebraic
interpretation. We spell it out with explicit coordinate and domain
conventions to distinguish the finite representation from the oscillator
representation.

\subsection{The quadratic oscillator realization}
On the Schwartz space $\mathcal S(\RR)$, let $D=d/dx$ and set
\begin{equation}\label{eq:oscillator-sl2}
\mathsf R=\frac{x^2}{2},\qquad
\mathsf L=\frac{D^2}{2},\qquad
\rho=xD+\frac12.
\end{equation}
The Heisenberg relation $[D,x]=1$ gives
\begin{equation}\label{eq:oscillator-sl2-brackets}
[\mathsf L,\mathsf R]=\rho,\qquad
[\mathsf L,\rho]=2\mathsf L,\qquad
[\rho,\mathsf R]=2\mathsf R.
\end{equation}
Thus this is an $\mathfrak{sl}(2)$ realization; a sign change of a
lowering generator gives another common convention.
The frequency-$\omega$ Hamiltonian and number operator are
\begin{equation}\label{eq:omega-oscillator}
\mathscr H_\omega=\frac12(-D^2+\omega^2x^2),\qquad
\mathcal N_\omega=\frac{\mathscr H_\omega}{\omega}-\frac12
=\omega\mathsf R-\omega^{-1}\mathsf L-\frac12.
\end{equation}
The natural Hilbert space for the wavefunctions and the unitary Fourier
transform is $L^2(\RR,dx)$. A Gaussian-weighted polynomial realization
requires conjugating the operators; it is not the same realization
with its measure changed silently.

\subsection{Squeezing is the exponential of $\rho$}
The closure of $\rho$ generates the unitary dilation group
\begin{equation}\label{eq:dilation}
V_r=e^{r\rho},\qquad (V_rf)(x)=e^{r/2}f(e^rx),\qquad r\in\RR.
\end{equation}
Direct calculation gives
\[
V_rxV_r^{-1}=e^rx,\qquad V_rDV_r^{-1}=e^{-r}D.
\]
For $e^{2r}=\omega$ this implies
\begin{equation}\label{eq:squeezed-conjugacy}
V_r\mathscr H_1V_r^{-1}=\frac{\mathscr H_\omega}{\omega},\qquad
\mathcal N_\omega=V_r\mathcal N_1V_r^{-1}.
\end{equation}
In particular,
\[
V_r(\pi^{-1/4}e^{-x^2/2})
=(\omega/\pi)^{1/4}e^{-\omega x^2/2}.
\]
Thus the changed width of our Perron profile is an actual dilation
in the oscillator representation, not only an analogy with squeezing.

\subsection{The limiting kernel as an algebraic semigroup}
Using \eqref{eq:omega-oscillator} in the variable $y$, the exact
Mehler operator can be written
\begin{equation}\label{eq:Lie-semigroup}
\boxed{
\mathcal G=c\,e^{\tau/2}
\exp\left[\tau(\alpha^{-1}\mathsf L-\alpha\mathsf R)\right],
\qquad \alpha=\sqrt8.
}
\end{equation}
Equivalently,
\[
\mathcal G=cV_{\frac12\log\alpha}
 e^{-\tau\mathcal N_1}V_{\frac12\log\alpha}^{-1}.
\]
For positive real $\tau$, this is a bounded imaginary-time oscillator
semigroup, rather than a unitary real-time rotation. Algebraically it
is the semigroup continuation of the same quadratic oscillator
construction used for fractional Fourier transformations.

With the positive-sign Fourier convention of the supplied notes,
\[
(\mathcal F_+f)(\xi)=\frac1{\sqrt{2\pi}}\int e^{i\xi x}f(x)\,dx,
\qquad
\mathcal F_+=e^{i\pi\mathcal N_1/2}.
\]
This contrasts with the negative-sign convention used in much of the
optics literature. The notes derive this identity from
$\psi_n=(x-D)^ne^{-x^2/2}$ and
$\mathscr H_1\psi_n=(n+1/2)\psi_n$.
In their generating-function calculation the intermediate exponent
is $-\xi^2/2+2i\xi t+t^2$; the factor $2$ in the mixed term is needed.
We use normalized Hermite wavefunctions and the unweighted $L^2$ space
when passing from the algebraic calculation to unitary operators.

\subsection{What this says, and does not say, about the finite algebra}
Equation \eqref{eq:T-enveloping} places the finite $T_N$ in the universal
enveloping algebra of \eqref{eq:finite-sl2}. Equation
\eqref{eq:Lie-semigroup} places the candidate limiting kernel in the
oscillator semigroup associated with \eqref{eq:oscillator-sl2-brackets}.
These are two different representations. We have not exhibited an
exact finite-$N$ Lie-group conjugacy between them, and $T_N$ is not
simply the standard normalized finite-oscillator generator.
The centered scaling of $H_N$, its Perron localization, and the
weighted kernel asymptotics are the analytic bridge that must be
controlled. The algebraic interpretation organizes that bridge but
does not replace its estimates.

The even and odd Hermite subspaces are each preserved by
$\mathsf R,\mathsf L,\rho$. This oscillator parity should not be
confused with the even/odd \emph{lattice-index} blocks of the finite
matrices. Each finite lattice block samples the entire real line in
the central scaling and can support both even and odd continuum modes.

Finally, the notes ask about an intertwiner with convolution.
Identifying a finite mode basis with the first $N+1$ Hermite modes
intertwines diagonal fractional phases on that subspace. It does not
by itself intertwine ordinary convolution with a finite Kravchuk
convolution, or identify the physical position operators. We leave
that stronger question separate from the spectral argument.
